\section*{الفصل \ref{ch:b3:total-synthesis} --- التخليق الكلي: الاستراتيجية والطرق الكلاسيكية}

\begin{solution}{exo:b3:total-synthesis:1}
بمردود 90\,\%: $0.9^{10} = 35$\,\% و$0.9^{20} = 12$\,\%. وبمردود 80\,\%: $0.8^{10} = 11$\,\% و$0.8^{20} = 1.2$\,\%.
\end{solution}

\begin{solution}{exo:b3:total-synthesis:2}
$5 + 3 + 1 + 4 = 13$ خطوة في المجموع؛ \omterm{def:b3:total-synthesis:architecture}{وأطول تسلسل خطي} $5 + 1 + 4 = 10$.
\end{solution}

\begin{solution}{exo:b3:total-synthesis:3}
4-(ثنائي ميثيل أمينو)بوتان-2-ون، $\ce{CH3COCH2CH2N(CH3)2}$: يُضاف إينول الأسيتون إلى أيون
الإيمينيوم $\ce{CH2=N+(CH3)2}$.
\end{solution}

\begin{solution}{exo:b3:total-synthesis:4}
افصل رابطتي الحلقة الأليليتين بالنسبة إلى C=C: البوتا-1،3-ديين والبيوت-3-إين-2-ون
(ميثيل فينيل كيتون)، اللذان يتفاعلان بالتسخين.
\end{solution}

\begin{solution}{exo:b3:total-synthesis:5}
الخطي: $0.9^{21} = 11$\,\%. المتقارب: $0.9^{11} = 31$\,\%، أي نحو ثلاثة أضعافه.
\end{solution}

\begin{solution}{exo:b3:total-synthesis:6}
A: $(7 + 1)/12 = 67$\,\%. B: $7/9 = 78$\,\%.
\end{solution}

\begin{solution}{exo:b3:total-synthesis:7}
دونهما: $0.9^8 = 43$\,\%. ومعهما، بأربع خطوات إضافية بمردود 95\,\%: $43\,\% \times 0.95^4 = 35$\,\%.
\end{solution}

\begin{solution}{exo:b3:total-synthesis:8}
الأولية: إيثر سيليلي ضخم (\textit{tert}-بوتيل ثنائي فينيل سيليل)، يُنزع بالفلوريد.
الثانوية: إيثر 4-ميثوكسي بنزيل، يُنزع بالأكسدة بالمركب DDQ. والكحول الثالثي،
المعاق، يمكن غالبًا أن يُترك حرًا، أو يُحمى في صورة إيثر سيليلي صغير
ينشطر بحمض معتدل. وكل ظرف يترك المجموعات الأخرى دون مساس.
\end{solution}

\begin{solution}{exo:b3:total-synthesis:9}
احتاجت التخاليق الكلية إلى عشرات الخطوات بمردودات إجمالية أدنى بكثير من 1\,\%؛ أما
مركب يحمل النظام الحلقي كاملًا ومعظم المراكز الفراغية، يُستخلص من
أوراق الطقسوس الأوروبي، وهو مصدر متجدد، فيُحوَّل في بضع خطوات.
\end{solution}

\begin{solution}{exo:b3:total-synthesis:10}
إضافة مايكل: يُضاف إينولات الديون (عند C2، بين مجموعتي الكربونيل) إلى
الكربون $\beta$ في ميثيل فينيل كيتون، فيعطي ثلاثي كيتون. الألدول: يهاجم إينولات
الميثيل كيتون كربونيلَ الحلقة داخل الجزيء، ويعطي نزع الماء
السيكلوهكسينون: كيتون فيلاند--ميشر، وهو إين ديون ثنائي الحلقة بمجموعة ميثيل
زاوية.
\end{solution}

\begin{solution}{exo:b3:total-synthesis:11}
$y^{m+1}/y^{2m+1} = y^{-m}$: بمردود 90\,\%، 1.7 من أجل $m = 5$ و4.9 من أجل $m = 15$. وكلما طال الطريق، ازداد
ربح التقارب.
\end{solution}

\begin{solution}{exo:b3:total-synthesis:12}
تنغلق كل حلقة والكاتيون والرابطة المضاعفة المهاجِمة في ترتيب شبيه
بالكرسي؛ والإضافة عبر كل رابطة مضاعفة من الجهتين المتقابلتين، وتأخذ المستبدلات
مواضع شبه استوائية، وهذا يجعل كل التحام حلقي \textit{trans}. والرابطة المضاعفة \textit{E}
تضع امتدادَي سلسلتها بحيث ينتج الالتحام \textit{trans}؛ أما الرابطة \textit{Z} فكانت ستعطي
الالتحام \textit{cis}: فهندسة البولي إين تُرمّز الكيمياء الفراغية
للناتج.
\end{solution}

\begin{solution}{pb:b3:total-synthesis:1}
\textbf{1.} $\ce{C4H6O2}$, $\ce{CH5N}$, $\ce{C5H6O5}$, $\ce{C8H13NO}$.

\textbf{2.} $\ce{C4H6O2 + CH5N + C5H6O5 -> C8H13NO + 2CO2 + 2H2O}$.

\textbf{3.} 86.0 و31.0 و\qty{146.0}{g/mol}؛ والتروبينون \qty{139.0}{g/mol}.

\textbf{4.} $139.0/(86.0 + 31.0 + 146.0) = 0.53$: أي 53\,\%.

\textbf{5.} مجموعتا $\ce{CH2}$ المركزيتان فيه، الواقعة كل منهما بين مجموعتي كربونيل، تتحولان بسهولة إلى إينول في
الماء؛ أما الأسيتون نفسه فضعيف المحبة للنواة جدًا، ومجموعتا الكربوكسيل
المنشِّطتان تغادران بعد ذلك.

\textbf{6.} تحتاج أيونات الإيمينيوم إلى أمين حر (والحمض الزائد يبرتنه)، ويحتاج
الإينول وتكوّن الإيمينيوم إلى قليل من الحمض: فيوازن pH قريب من التعادل بينهما.

\textbf{7.} يُضاف $\ce{CH3NH2}$ إلى إحدى مجموعتي الألدهيد؛ ويعطي فقد الماء أيون الإيمينيوم
$\ce{R-CH=N+H-CH3}$.

\textbf{8.} يهاجم كربون إينولي في حمض الأسيتون ثنائي الكربوكسيليك كربون الإيمينيوم، فيكوّن
رابطة C--C وأمينًا ثانويًا.

\textbf{9.} يتكاثف الأمين الثانوي مع مجموعة الألدهيد الثانية في السلسلة
نفسها، فيعطي أيون إيمينيوم حلقيًا (خماسيًا).

\textbf{10.} يهاجمه الإينول الواقع على الجانب الآخر من الكيتون، فيغلق
الحلقة السداسية والهيكل الثنائي الحلقة.

\textbf{11.} كل منهما حمض $\beta$-كيتوني، يفقد $\ce{CO2}$ عبر حالة انتقالية حلقية سداسية،
فيعطي إينول الكيتون.

\textbf{12.} رابطتا C--C ورابطتا C--N.

\textbf{13.} $0.75^{14} = 0.018$: أي 1.8\,\%.

\textbf{14.} $42/1.78 = 24$.

\textbf{15.} 100\,\%: خطوة واحدة، كلها بناء روابط.

\textbf{16.} عملية واحدة، والماء مذيبًا، ومكوّنات رخيصة، ولا مجموعات حامية 
وتكاد لا توجد نفايات غير ثنائي أكسيد الكربون والماء.

\textbf{17.} يبني النبات الهيكل التروباني من أيون إيمينيوم حلقي مشتق من حمض
أميني ومن وحدة مشتقة من الأسيتات، عبر كيمياء من نوع مانيخ
نفسه.

\textbf{18.} التفاعلات الشلالية (والإضافات الحلقية المكوّنة للحلقات مثل تفاعل
ديلز--ألدر).

\textbf{19.} لا: \omterm{def:b3:point-groups:mirror}{فمستوي مرآة} يمر بالذرة N ومجموعة الميثيل فيها وC3 وأكسجين
الكربونيل.

\textbf{20.} لكل منهما أربع مجموعات مختلفة، لكنهما صورتان مرآويتان إحداهما للأخرى في
مستوي التناظر الخاص بالجزيء: مركب من نمط ميزو.

\textbf{21.} متماكبان لامرآويان، يختلفان في اتجاه المجموعة OH على C3 نسبةً إلى
الجسر النيتروجيني.

\textbf{22.} لا: \omterm{def:b3:point-groups:mirror}{فمستوي المرآة} محفوظ في كليهما.

\textbf{23.} وجها مجموعة الكربونيل غير متكافئين: أحدهما يواجه
الجسر النيتروجيني، والآخر الجسر ذا ذرتي الكربون، فيعيقان على نحو مختلف.

\textbf{24.} يعطي الطريق في وعاء واحد نحو 24 ضعف \omterm{def:b3:total-synthesis:architecture}{المردود الإجمالي} للطريق الخطي ذي
الأربع عشرة خطوة.
\end{solution}
